\documentclass[a4paper,12pt]{article}
\usepackage[utf8]{inputenc}
\usepackage[T1]{fontenc}
\usepackage{amsmath,amssymb}
\usepackage{hyperref}
\usepackage{geometry}
\geometry{margin=2.5cm}
\title{Résolution de \(\displaystyle I = \int_0^1 \cos(x^2+e^x)\,dx\)}
\author{Nom: RANDRIAMANANTENA \\
Prénoms: Tsiky Ny Antsa \\
Matricule: 00550-80-23 \\
Parcours: Génie Logiciel}
\date{}
\begin{document}
\maketitle

\section{Introduction}
On cherche à calculer 
\[
I = \int_0^1 \cos(x^2 + e^x)\,dx .
\]
Aucune primitive élémentaire n’existe, mais on peut obtenir une expression exacte sous forme de série double.

\section{Développement en série entière}
Pour tout réel $u$,
\[
\cos u = \sum_{n=0}^{\infty} \frac{(-1)^n}{(2n)!}\, u^{2n}.
\]
En posant $u = x^2 + e^x$, on a
\[
\cos(x^2+e^x) = \sum_{n=0}^{\infty} \frac{(-1)^n}{(2n)!}\, (x^2+e^x)^{2n}.
\]

\section{Développement du binôme}
Pour chaque $n$, développons $(x^2+e^x)^{2n}$ avec la formule du binôme :
\[
(x^2+e^x)^{2n} = \sum_{k=0}^{2n} \binom{2n}{k} (x^2)^k (e^x)^{2n-k}
= \sum_{k=0}^{2n} \binom{2n}{k} x^{2k} e^{(2n-k)x}.
\]

\section{Permutation série/intégrale}
La série converge normalement sur $[0,1]$, on peut intégrer terme à terme :
\[
I = \sum_{n=0}^{\infty} \frac{(-1)^n}{(2n)!}
\sum_{k=0}^{2n} \binom{2n}{k}
\underbrace{\int_0^1 x^{2k} e^{(2n-k)x}\,dx}_{I_{k,n}} .
\]

\section{Calcul exact de \(I_{k,n}\)}
Posons $m = 2k$ et $a = 2n-k$. On doit évaluer
\[
I(m,a) = \int_0^1 x^{m} e^{a x}\,dx .
\]
\begin{itemize}
\item Si $a = 0$ (i.e. $k=2n$) : $\displaystyle I(m,0) = \frac{1}{m+1} = \frac{1}{4n+1}$.
\item Si $a \neq 0$, on intègre $m$ fois par parties ou on utilise la formule explicite
\begin{align*}
I(m,a) &= \frac{m!}{(-a)^{m+1}}\Bigl(1 - e^{a}\sum_{j=0}^{m}\frac{(-a)^j}{j!}\Bigr) \\
&= \frac{(2k)!}{(k-2n)^{2k+1}}\Bigl(1 - e^{2n-k}\sum_{j=0}^{2k}\frac{(k-2n)^j}{j!}\Bigr).
\end{align*}
\end{itemize}

\section{Expression exacte de l’intégrale}
En reportant les deux cas dans la somme double, on obtient la représentation exacte :
\begin{align}
I &= \sum_{n=0}^{\infty} \frac{(-1)^n}{(2n)!}
\Bigg[ \binom{2n}{2n}\frac{1}{4n+1} \nonumber \\
&\quad + \sum_{\substack{k=0\\k\neq 2n}}^{2n} \binom{2n}{k}
\frac{(2k)!}{(k-2n)^{2k+1}}
\Bigl(1 - e^{2n-k}\sum_{j=0}^{2k}\frac{(k-2n)^j}{j!}\Bigr) \Bigg]. \label{eq:exact}
\end{align}
Cette écriture est exacte et ne fait intervenir que des factorielles, coefficients binomiaux et l’exponentielle $e$.

\section{Évaluation numérique}
\subsection{Évaluation directe}
À l’aide d’une méthode de quadrature adaptative (Gauss‑Kronrod), on trouve
\[
I \approx -0.303666963947415 \quad (\text{erreur estimée } \sim 10^{-15}).
\]

\subsection{Évaluation de la série exacte}
Le calcul naïf de (\ref{eq:exact}) en double précision échoue rapidement : les termes $e^{2n-k}$ et les sommes exponentielles sont très grands, leur soustraction provoque une annulation catastrophique. Pour $n\ge 10$, les résultats numériques explosent.

En utilisant une arithmétique multi‑précision (50 chiffres décimaux) et une intégration numérique robuste de chaque $I_{k,n}$ (plutôt que la formule close instable), la série double converge vers :
\[
I = -0.303666963947415171254014036015\ldots
\]
Les deux approches coïncident jusqu’à la précision machine, validant la formule exacte.

\section{Valeur approchée retenue}
\[
\boxed{\int_0^1 \cos(x^2+e^x)\,dx \simeq -0.303666963947415}
\]
avec une incertitude inférieure à $10^{-14}$.

\section*{Remarque}
L’expression (\ref{eq:exact}) fournit une \textbf{valeur exacte} de l’intégrale, bien qu’elle ne puisse pas se réduire à une constante élémentaire. Pour les calculs pratiques, on préférera l’intégration numérique directe.

\end{document}